\section{半棋盘双射与第二种证明}\label{sec:half}
本节从局部覆盖重新构造计数，不使用基本块分类或回文公式。
为便于后续推广，先允许任意整数高度 $m\ge1$。

\subsection{带末端半块的部分铺法}
对 $j\ge1$ 和 $S\subseteq\{1,\ldots,m\}$，令 $\mathcal P_m(j,S)$ 为
$m\times j$ 棋盘的部分铺法：除末列行号属于 $S$ 的格外，所有格都由内部多米诺覆盖；
这些未覆盖格标记为向右伸出的半块，准备与右邻格配成水平砖。
$j=0$ 时只有 $S=\varnothing$ 的空对象。记 $q_j(S)=|\mathcal P_m(j,S)|$。

令 $T_m$ 的行列按所有行子集排列。定义
\begin{equation}\label{eq:Tm}
 T_m[p,q]=
 \begin{cases}
 1,&p\cap q=\varnothing,\text{ 且 }[m]\setminus(p\cup q)
       \text{ 的每个连续段均为偶长},\\
 0,&\text{否则}.
 \end{cases}
\end{equation}
这里 $[m]=\{1,\ldots,m\}$，空补集满足条件。
每段偶长剩余行的竖直配对唯一，因为端点只能与相邻格配对，
剥去此对后继续。奇长段无法配满。
因此 \eqref{eq:Tm} 恰好刻画加一列的合法转移。
放入指定水平砖和唯一竖直配对给出逆构造，故
\begin{equation}\label{eq:profile}
 q_j^{\trans}=\evec^{\trans}T_m^j,\qquad
 t_m(j)=\evec^{\trans}T_m^j\evec,
\end{equation}
其中 $\evec$ 是空状态的单位列向量，$t_m(j)$ 为完整 $m\times j$ 铺法数。
辅助定义 $T_0=(1)$、$t_0(j)=1$。

\subsection{对称固定铺法的边界条件}
记 $J(S)=\{m+1-r:r\in S\}$，$\one$ 为全 1 列向量，
$\chi_m(S)=1$ 当且仅当 $J(S)=S$。
用 $v_{m,n},r_{m,n}$ 表示 $m\times2n$ 的左右反射、半周旋转固定数。

\begin{lemma}[中线补全双射]\label{lem:half}
对 $m\ge1$ 和 $n\ge0$，
\begin{equation}\label{eq:half-fixed}
 v_{m,n}=\evec^{\trans}T_m^n\one,\qquad
 r_{m,n}=\evec^{\trans}T_m^n\chi_m.
\end{equation}
\end{lemma}
\begin{proof}
先设 $n\ge1$。任取 $P\in\mathcal P_m(n,S)$，
将内部砖左右反射到右半棋盘，然后对每个 $r\in S$
插入连接 $(r,n)$ 与 $(r,n+1)$ 的水平砖。
两半内部砖不相交，跨中线砖恰好补齐剩余格，
得到唯一的左右反射固定铺法。
反向限制到左半，切断跨中线砖，唯一恢复 $P,S$，
因此所有 $S$ 都可取，并有第一式。

对半周旋转，将左半内部砖旋转到右半后，右侧待补行集合为 $J(S)$。
跨中线砖必须水平，故能补全的充要条件是 $J(S)=S$。
条件满足时逐行插入跨中线砖，所得砖集合在半周旋转下保持不变；
限制到左半是逆操作，得到第二式。
条件的必要性也来自跨中线第 $r$ 行的砖必须有第 $m+1-r$ 行的旋转像。
$n=0$ 时只有空对象，两式均为 1。
\end{proof}

\begin{lemma}[奇数高度的上下反射]\label{lem:h-odd}
若 $m=2k+1$，$k\ge0$，则
\begin{equation}\label{eq:h-odd}
 h_{m,n}=t_k(2n).
\end{equation}
\end{lemma}
\begin{proof}
在上下反射固定铺法中，接触中央行的竖直砖会被映到另一块竖直砖，
两者在中央行重叠，因此不存在这样的砖。
中央行只能以唯一方式水平铺满。其上 $k$ 行必须自行铺满，
下方由其反射唯一确定。任意上方完整铺法都给出一个固定铺法，
限制回上方为逆操作。$k=0$ 与 $n=0$ 由空对象约定覆盖。
\end{proof}

\subsection{三行的两维递推}
对 $m=3$，由 \eqref{eq:Tm} 从空状态可达的状态恰为
$\varnothing,\{1\},\{3\},\{1,2\},\{2,3\},\{1,2,3\}$。
规则给出空集的输出为两个端点单集和全体集合；
$\{1\}$ 的输出为 $\varnothing,\{2,3\}$，
$\{3\}$ 的输出为 $\varnothing,\{1,2\}$，
其余三个状态的输出分别被迫为 $\{3\},\{1\},\varnothing$。
下面将状态支持与计数同时归纳。记
\[
 E=\{\varnothing,\{1,2\},\{2,3\}\},\qquad
 O=\{\{1\},\{3\},\{1,2,3\}\}.
\]
完整转移表表明，$E$ 的全部出口在 $O$，$O$ 的全部出口在 $E$，
没有通向剩余两个状态的出口。这只表示允许的转移交替，
不表示两集合间任意一对状态都有转移。

\paragraph{归纳命题与基例。}
对每个 $j\ge0$，$q_{2j}$ 的支持包含于 $E$，且
\[
 q_{2j}(\varnothing)=\alpha_j,\qquad
 q_{2j}(\{1,2\})=q_{2j}(\{2,3\})=\beta_j.
\]
这里 $\beta_j$ 计数一个指定状态，两种非空状态合计为 $2\beta_j$。
支持包含于 $E$ 不要求其中每个状态都有正计数。
零宽只有空对象，所以 $(\alpha_0,\beta_0)=(1,0)$，
两个非空状态的计数相等且均为零。

\paragraph{增加第一列。}
假设 $2j$ 列时命题成立。按目标状态的全部来源求和，得到
\[
\begin{aligned}
q_{2j+1}(\{1\})
 &=q_{2j}(\varnothing)+q_{2j}(\{2,3\})=\alpha_j+\beta_j,\\
q_{2j+1}(\{3\})
 &=q_{2j}(\varnothing)+q_{2j}(\{1,2\})=\alpha_j+\beta_j,\\
q_{2j+1}(\{1,2,3\})&=q_{2j}(\varnothing)=\alpha_j.
\end{aligned}
\]
其余状态计数为零，因为 $E$ 的出口已穷尽。
每条合法单列转移的局部补法唯一；反过来，移除新增列并切断跨越其左边界的水平砖，
可恢复唯一旧对象和输入状态。因此各来源不重不漏，求和时不需要额外乘数。

\paragraph{增加第二列。}
$O$ 中三个状态都能转入空集；只有 $\{3\}$ 能转入 $\{1,2\}$，
只有 $\{1\}$ 能转入 $\{2,3\}$。于是
\[
\begin{aligned}
q_{2j+2}(\varnothing)
 &=q_{2j+1}(\{1\})+q_{2j+1}(\{3\})+q_{2j+1}(\{1,2,3\})\\
 &=(\alpha_j+\beta_j)+(\alpha_j+\beta_j)+\alpha_j
   =3\alpha_j+2\beta_j,\\
q_{2j+2}(\{1,2\})&=q_{2j+1}(\{3\})=\alpha_j+\beta_j,\\
q_{2j+2}(\{2,3\})&=q_{2j+1}(\{1\})=\alpha_j+\beta_j.
\end{aligned}
\]
因 $O$ 的全部出口在 $E$，其余状态仍为零。
这同时证明下一偶数列的支持和两个非空状态的计数相等性，
从而完成对全部 $j\ge0$ 的归纳。最后写成列向量递推：
\begin{equation}\label{eq:two-dimensional}
 \binom{\alpha_{j+1}}{\beta_{j+1}}
 =\begin{pmatrix}3&2\\1&1\end{pmatrix}
 \binom{\alpha_j}{\beta_j},\qquad
 (\alpha_0,\beta_0)=(1,0).
\end{equation}

\paragraph{四个矩阵系数的含义。}
第一行计数到达空集的对象；第二行只计数到达 $\{1,2\}$ 的对象。
\begin{center}
\small
\begin{tabular}{cp{0.78\linewidth}}
\toprule
系数 & 两列接续的计数理由\\
\midrule
$3$ & 从空集经 $\{1\}$、$\{3\}$ 或 $\{1,2,3\}$ 回空集，共三条路径。\\
$2$ & 从 $\{1,2\}$ 和 $\{2,3\}$ 各有一条回空集的路径；两组来源各有 $\beta_j$ 个对象。\\
$1$ & 从空集经 $\{3\}$ 到 $\{1,2\}$，有且仅有一条路径。\\
$1$ & 从 $\{1,2\}$ 经 $\{3\}$ 回到自身；另一非空状态不能在两步到达此目标。\\
\bottomrule
\end{tabular}
\end{center}
特别地，系数 $2$ 不是说每个固定非空状态对象有两种回空接法。
例如 $(\alpha_1,\beta_1)=(3,1)$ 给出
$\alpha_2=(3+1)+(3+1)+3=11$、$\beta_2=3+1=4$。
这个实例用于解释已证明的归纳，不作为一般结论的依据。

\begin{proposition}\label{prop:parity}
设 $d_j=t_3(2j)$，则
\begin{equation}\label{eq:d-rec}
 d_0=1,\quad d_1=3,\quad d_{j+2}=4d_{j+1}-d_j\quad(j\ge0).
\end{equation}
三行固定点数满足
\begin{equation}\label{eq:parity-fixed}
 h_n=1,\quad r_{2j}=r_{2j+1}=d_j,\quad
 v_{2j}=d_{j+1}-2d_j,\quad v_{2j+1}=d_{j+1}.
\end{equation}
因此
\begin{equation}\label{eq:parity-orbits}
 a(2j)=\frac{d_{2j}+1+d_{j+1}-d_j}{4},\qquad
 a(2j+1)=\frac{d_{2j+1}+1+d_{j+1}+d_j}{4}.
\end{equation}
\end{proposition}
\begin{proof}
由 \eqref{eq:profile}，$d_j=\alpha_j$。
式 \eqref{eq:two-dimensional} 中的矩阵满足 $M^2-4M+I=0$，
从而得 \eqref{eq:d-rec}。
偶数步可达的 $J$-不变集合只有空集，
奇数步可达的 $J$-不变集合只有全三行集合，故引理 \ref{lem:half} 给出 $r$ 的两式。
对所有状态求和，偶数步得
$\alpha_j+2\beta_j=d_{j+1}-2d_j$，
奇数步得 $3\alpha_j+2\beta_j=d_{j+1}$。
引理 \ref{lem:h-odd} 给出 $h_n=t_1(2n)=1$。
最后代入轨道平均公式即得 \eqref{eq:parity-orbits}。
\end{proof}

\begin{proof}[定理 \ref{thm:main} 的第二种证明]
递推 \eqref{eq:d-rec} 给出 $D(x)=(1-x)/(1-4x+x^2)$。
对 \eqref{eq:parity-fixed} 分别求和可得
\[
 H(x)=\frac1{1-x},\quad R(x)=(1+x)D(x^2),\quad
 V(x)=(1+x)\frac{D(x^2)-1}{x^2}-2D(x^2).
\]
这里 $D(x^2)-1$ 可唯一写为 $x^2 E(x)$，
其中 $E(x)=\sum_{j\ge0}d_{j+1}x^{2j}$；上式的商表示这一系数移位，
并未将 $x$ 或 $x^2$ 当作形式幂级数环中的单位。
化简 $V$ 得到 \eqref{eq:V}，再用轨道平均公式即得
\eqref{eq:average-rational}，其通分和系数比较已在前文完整写出。
本次推导的组合前提仅为半棋盘双射和单列覆盖规则。
\end{proof}
